\secnumbersection{DEFINICIÓN DEL PROBLEMA}

\textcolor{pink}{Se debe definir el problema, es importante no confundir definir el problema con describir la solución. Por ejemplo: ``diseñar una arquitectura e implementar una plataforma ...'' es una solución, no un problema.}

\textcolor{pink}{Algunos elementos que podrían ir en este capítulo son (no es necesario que vayan todos):
\begin{itemize}
    \item Breve descripción del contexto donde se realizará la memoria (organización, línea dentro de la Informática en la que se basa, etc.)
    \item ¿Qué y cómo se realiza actualmente la situación que mejorarás con tu memoria?
    \item ¿Qué actores o usuarios están involucrados?
    \item ¿Qué dificultades tienen esos actores actualmente? ¿cuántos son? (ideal si se pueden poner estadísticas para así saber si existe un mercado razonable para la solución que propondrás en tu memoria, en el fondo saber cuántas personas u organizaciones tienen el mismo problema que estás definiendo)
    \item ¿Qué podría pasar si en el corto o mediano plazo no se solucionan esas dificultades (¿es decir, si no se hiciera tu memoria, qué pasaría?; en el fondo justificar por qué conviene hacer tu memoria, ¿cuál es la motivación o interés de hacerla?).
    \item ¿Qué competencia existe actualmente? (a lo mejor ya existe una solución al problema, pero por qué no sirve, o por qué tu solución sería mejor, también se puede enfocar a si este problema existe en otras realidades y cómo ha sido solucionado allí).
    \item Precisar los objetivos y alcances de la memoria (o solución al problema).
\end{itemize}}

\textcolor{pink}{En este capítulo, de ser necesario puede usar referencias bibliográficas (velar porque sean recientes), una cita de ejemplo \cite{schwab2002cure} y otras más \cite{georget1994study,beaumont1990patient}.}

\textcolor{pink}{Recuerde poner notas al pie de página que sean explicativas \footnote{Este es un ejemplo de una nota al pie de página. Puede indicar alguna URL, definiciones, aclarar alguna información pertinente del texto, citar algunas referencias, etc..}.}

\subsection{PROBLEMÁTICA SOCIAL}

Las personas con discapacidad visual enfrentan barreras persistentes al interactuar con entornos digitales. La Organización Mundial de la Salud reporta que al menos 2.200 millones de personas en el mundo tienen algún grado de deficiencia visual \cite{who2026blindness}, pero esa cifra agrupa mayoritariamente casos leves y corregibles con lentes. Al filtrar por los niveles de discapacidad relevantes para esta memoria —baja visión, discapacidad visual severa y ceguera total—, las estimaciones del Global Burden of Disease para el año 2020 son de 43,3 millones de personas ciegas y 295 millones con discapacidad visual moderada a severa, es decir, un orden de 338 millones de personas a nivel mundial \cite{bourne2021trends}.

Sin embargo, la magnitud de esta cifra no es lo relevante, sino la pertinencia ética de integrar a un grupo de usuarios históricamente segregado del entorno tecnológico general, en la medida en que se ve en la obligación de depender de herramientas diseñadas exclusivamente ``para ciegos'' en lugar de acceder al mismo contenido que el resto de los usuarios. Evitar segregar o estigmatizar a un grupo de usuarios mediante soluciones separadas ha sido, de hecho, uno de los principios formales del diseño universal —el principio de uso equitativo— en los últimos 30 años \cite{centerforuniversaldesign1997}.

La consecuencia de no resolver estas barreras es la exclusión en las dimensiones social, laboral y educativa. La imposibilidad de acceder de forma autónoma a contenido en línea limita:

\begin{itemize}
    \item \textbf{Social:} la participación social cotidiana.
    \item \textbf{Laboral:} el desempeño y las oportunidades laborales.
    \item \textbf{Educativo:} el acceso a instancias educativas, que involucran a un grupo más acotado.
\end{itemize}

A modo de ejemplo, hoy no es posible enterarse por cuenta propia de qué muestra una fotografía que comparte un familiar o un amigo, ni completar sin ayuda una compra o un trámite en línea cuando de por medio se exige una verificación visual \cite{webaim2024screenreadersurvey10}. Tampoco es posible darse cuenta a tiempo de un aviso que aparece de forma inesperada en la pantalla, lo que puede significar perder una oportunidad o un plazo importante. En el trabajo, muchas veces es necesario pedirle a un compañero que describa una cifra o una imagen antes de poder tomar una decisión, y postular a un empleo puede volverse una barrera en sí misma cuando el proceso exige una verificación que solo puede completarse viendo la pantalla. Algo similar ocurre en lo educativo, donde acceder a un material de estudio compartido como imagen deja a la persona a la espera de que alguien más se lo describa antes de poder seguir aprendiendo al mismo ritmo que el resto.

\textcolor{red}{fuente pendiente: confirmar y referenciar estudio(s) específicos sobre el impacto práctico de estas barreras en el ámbito laboral y educativo.}

\subsubsection{Búsqueda de antecedentes locales}

El punto de partida se basó en que esta problemática social afecta de manera particular a las personas ciegas al navegar la web. Sin embargo, al revisar el estado del arte no se encontraron estudios específicos que caracterizaran con suficiente detalle esa necesidad en este grupo, sobre todo porque ya existen soluciones tecnológicas que resuelven parte de esto, como por ejemplo los lectores de pantalla \footnote{Programas que traducen el contenido de una página web a voz sintetizada o a un dispositivo braille; se profundiza en su funcionamiento en el Capítulo 2.}. Por ello, la problemática social se acotó hacia el contexto específico de personas con discapacidad visual que usan lector de pantalla pero aún presentan dificultades para utilizar la plenitud de la información que ofrece un sitio web al momento completar una tarea, y se diseñó y aplicó una encuesta propia con el fin de validar empíricamente la existencia y la naturaleza de la necesidad antes de proponer una solución.

\subsubsection{La encuesta (trabajo realizado)}

Se elaboró y distribuyó el instrumento ``Encuesta sobre tecnologías para personas ciegas'' mediante Google Forms, enviado por correo electrónico a miembros de la Fundación FUNDALURP (fundación de perros guía para la rehabilitación de personas en situación de discapacidad visual). La encuesta se aplicó entre febrero y marzo de 2026 y se obtuvieron 9 respondientes. El detalle completo del instrumento —su estructura, tipos de pregunta y cómo se relacionan con la problemática— se presenta en el Anexo \ref{anexo:encuesta}.

El perfil de los respondientes muestra como condición dominante la retinitis pigmentaria, con niveles de discapacidad que van desde severo hasta ceguera total, y un manejo computacional autoreportado mayoritariamente entre básico e intermedio.

Respecto de las dificultades reportadas al navegar la web con lectores de pantalla, las más frecuentes, en orden, fueron: elementos estrictamente visuales que no pueden interpretarse (imágenes, gráficos); pop-ups o notificaciones que no se leen; campos de escritura sin etiqueta; la obligación de usar el ratón para ciertas interacciones; y los captchas visuales. \textcolor{red}{pendiente: agregar conteo o porcentaje exacto de respondientes que reportó cada dificultad.}

En cuanto a las funcionalidades más demandadas, se destacan tres, en orden de frecuencia: describir e interpretar imágenes; la navegación por voz hacia un objetivo concreto; y la lectura de notificaciones o etiquetas dinámicas. \textcolor{red}{pendiente: agregar el conteo o porcentaje exacto de respondientes que solicitó cada funcionalidad.}

\textcolor{red}{gráficos: distribución del perfil de respondientes (condición, nivel de discapacidad, nivel computacional) y frecuencia de dificultades y demandas reportadas.}

\subsection{PROBLEMÁTICA TÉCNICA}

Actualmente, las personas con discapacidad visual navegan la web apoyándose en lectores de pantalla, programas como JAWS, NVDA o VoiceOver que traducen el contenido de una página a voz sintetizada o a un dispositivo braille. El DOM (Document Object Model) es la estructura de nodos que el navegador construye a partir del código HTML de la página; a partir de ella, el navegador deriva un árbol de accesibilidad —una representación paralela que expone únicamente los roles, nombres accesibles, estados y relaciones relevantes para una tecnología asistiva— y lo pone a disposición a través de la API de accesibilidad del sistema operativo \cite{mdn2024accessibilitytree}. Es esa API la que el lector de pantalla consulta para construir lo que finalmente comunica al usuario. También existe la posibilidad de que el usuario navegue por una página utilizando únicamente el teclado, lo cual es posible gracias al atributo \texttt{tabindex} y al orden del DOM. El recorrido secuencial mediante la tecla Tab lo controla el propio navegador para la navegación por teclado, en el orden del DOM o según el atributo \texttt{tabindex}, y no es exclusivo de los lectores de pantalla ni de los usuarios con discapacidad visual.

Buena parte de la información que un lector de pantalla puede comunicar depende de las etiquetas ARIA (Accessible Rich Internet Applications): un conjunto de atributos que se agregan a los elementos del DOM para declarar roles, estados y relaciones que no tienen una etiqueta HTML nativa que los describa \cite{mdn2024aria}. Por ejemplo, una notificación emergente puede marcarse como una región viva (\texttt{aria-live}) para que el lector la anuncie apenas aparece. Cuando un sitio no declara estas etiquetas, el lector de pantalla no tiene forma de identificar ese contenido dinámico, aunque técnicamente sí forme parte del DOM.

El límite de este enfoque está, entonces, en que depende por completo de que la página declare la semántica correcta —HTML nativo y ARIA— para cada elemento. Los lectores de pantalla no interpretan información estrictamente visual: no describen imágenes, gráficos ni captchas, y solo leen de forma confiable las notificaciones o pop-ups que fueron etiquetados correctamente por quien construyó el sitio. Tampoco comprenden la estructura espacial de la interfaz (qué elemento está dónde y con qué propósito) más allá del orden en que fueron declarados. El usuario recibe así una representación incompleta de la página, que depende de decisiones de accesibilidad ajenas a su control.

Respecto de la competencia y soluciones existentes, las alternativas actuales resultan insuficientes. Los lectores de pantalla, por diseño, no cubren la interpretación visual; y las soluciones basadas en IA para descripción de imágenes no están integradas al flujo de navegación ni fueron pensadas específicamente para personas ciegas, por lo que no resuelven el problema de forma completa.

\textcolor{red}{tabla: comparación de soluciones existentes (lectores de pantalla, herramientas de descripción de imágenes por IA etc) frente a las capacidades requeridas, para justificar la brecha que cubre la propuesta.}

\textcolor{red}{fuente pendiente: referencias recientes sobre lectores de pantalla, accesibilidad web (p. ej. WCAG) y soluciones de IA para descripción de imágenes.}

\subsection{OBJETIVOS Y ALCANCES}

Considerando la relevancia de que las personas con discapacidad visual puedan acceder de forma equitativa a las tecnologías de la información, se plantean los siguientes objetivos para el desarrollo del asistente de accesibilidad web que en este trabajo se denomina Relator.

\textbf{Objetivo general.} Desarrollar el Relator, un sistema basado en inteligencia artificial capaz de analizar e interpretar la información visual de sitios web a partir de capturas de pantalla, utilizando modelos de lenguaje visual (VLM), que se integre y coexista con el lector de pantalla del usuario, con el fin de mejorar la accesibilidad de las personas con discapacidad visual.

\textbf{Objetivos específicos.}

\begin{itemize}
    \item[(a)] Implementar un prototipo funcional del asistente (Relator) capaz de interpretar información visual y entregar descripciones auditivas.
    \item[(b)] Definir instrucciones de navegación que orienten al modelo en la interpretación de los elementos web.
    \item[(c)] Empaquetar e implementar la solución en un entorno accesible y adaptable.
    \item[(d)] Validar la solución mediante pruebas con usuarios reales con discapacidad visual.
\end{itemize}

\textbf{Alcances.} \\Se espera que la aplicación funcione sobre navegadores basados en Chromium (Google Chrome, Microsoft Edge, Opera, entre otros) cubriendo la mayoría de los casos de uso. \textcolor{red}{(Incluir cifras de uso de navegadores)} El idioma de interacción contemplado es el español, aunque se espera que el modelo VLM seleccionado soporte de forma nativa otros idiomas \textcolor{red}{(Esto se confirmará en Capitulo 3. Propuesta de Solución)}.
